% Separate post-CP15 research. Not assembled pending independent audits.
\begin{proposition}[Proved: linear-size Eisenstein charts and pointed isomorphism witnesses]
\label{D:local-Eisenstein-linear-atlas}
Let $k$ be an algebraically closed field of arbitrary characteristic.
For integers $e\ge2$ and $d\ge e-1$, put
\[
                  N=\left\lfloor\frac{2d}{e}\right\rfloor+1.
                                                               \tag{D.EA.1}
\]
There is an explicitly defined affine $k$-scheme $\mathcal E_{e,d}$
with the following properties; it is allowed to be empty.
\begin{enumerate}
\item Its coordinates are at most $e(N-1)\le2d$ polynomial
coefficients and at most two auxiliary inverses. Its equations are
linear coefficient vanishings and degree-two inverse equations.
For every $k$-algebra $R$, a point determines a monic Eisenstein polynomial
\[
 f(X)=X^e+\sum_{i=1}^{e}a_i(t)X^{e-i},\qquad
 a_i(t)=\sum_{j=1}^{N-1}a_{ij}t^j.                       \tag{D.EA.2}
\]
The algebra $B_f=R[[t]][X]/(f)$ is finite free of rank $e$ over
$R[[t]]$, is canonically a pointed formal disc $R[[\alpha]]$ over $R$,
and satisfies
\[
 t=\alpha^e u,\qquad f'(\alpha)=\alpha^d v,
 \qquad u,v\in R[[\alpha]]^\times.                       \tag{D.EA.3}
\]
In particular $B_f[1/t]$ is finite etale over $R((t))$, and its
discriminant ideal in $R[[t]]$ is $(t^d)$.
\item Every separable extension of $k((t))$ of degree $e$ and
different exponent $d$, with the valuation on its integer ring
normalized to have value group $\mathbb Z$, is represented by a
$k$-point of this chart. The same assertion holds after every
algebraically closed extension of $k$.
\item Over $\mathcal E_{e,d}\times\mathcal E_{e,d}$ there is an affine
finite-presentation scheme $\mathcal W_{e,d}$ using, in total,
at most $7d+6$ coordinates and equations of degree at most $e^2+1$.
For every $k$-algebra $R$ and every pair $(f,g)$ over $R$, its fibre
has an $R$-point if and only if there exists a continuous
$R[[t]]$-algebra isomorphism $B_g\simeq B_f$ preserving the
specified augmentations to $R$. Thus the witness is effective on
arbitrary, including nonreduced, coefficient rings.
\end{enumerate}
Here a pointed formal disc means $R[[\alpha]]$ with its augmentation
$\alpha\mapsto0$ and its $\alpha$-adic topology. The assertion about
families is deliberately about augmentation-preserving isomorphisms;
it makes no assertion that arbitrary nilpotent families of unpointed
Laurent-series algebras have this presentation. At geometric points
the distinction does not affect the represented local field extensions.
The case $e=1,d=0$ is the single trivial extension and needs no parameters.
\end{proposition}
\begin{proof}
We give the equations and the lifting argument, so that the precision
and the family scope can be checked directly. The usual Eisenstein
coefficient spaces and their fixed-discriminant strata are treated in
\citet[Section 4, particularly Proposition 4.4, Corollary 4.6 and
Lemma 4.7]{Yasuda2025EisensteinB}. That source also distinguishes the
relative etaleness condition from merely having separable field fibres.
Our smaller coefficient cutoff follows from the explicit Newton
argument below; we do not use a local-field counting formula over a
finite residue field.

\emph{The coefficient stratum.}
Always require $a_{e1}$ to be a unit. If the integer $e$ is nonzero
in $k$, the chart is empty unless $d=e-1$. In the latter case no
additional derivative conditions are imposed. Suppose instead that
$e=0$ in $k$. For every $1\le i<e$ for which $e-i$ is nonzero in $k$,
assign the coefficient $a_{ij}$ the derivative weight
\[
                         ej+e-i-1.                    \tag{D.EA.4}
\]
Set it equal to zero whenever this weight is less than $d$.
There must be a pair $(i_*,j_*)$ of these indices with weight exactly
$d$; if there is none the chart is empty. There is at most one such
pair, because the residues $e-i-1$ modulo $e$ are distinct. Require
$a_{i_*j_*}$ to be a unit. These are the asserted equations, with
one inverse for each of the two required units. The required index
$j_*$, when it exists, is less than $N$: its weight is at least
$ej_*$, whereas $eN>2d\ge d$.

The same valuation argument includes the characteristic-zero and
tame cases. In fact the derivative of the leading monomial has
weight $e-1$ when it is nonzero. All other nonzero derivative terms
have weights (D.EA.4), with distinct residues modulo $e$. Thus the
least weight cannot cancel, and its coefficient is a unit under
exactly the stated conditions. In particular these equations encode
the specified different exponent on every field fibre without
introducing resultant equations of growing length.

\emph{Formal discs over arbitrary coefficient rings.}
Write $f(X)=X^e+tU(t,X)$. The constant term $U(0,0)=a_{e1}$ is
invertible. The formal implicit-function recursion, using this unit,
produces a unique series $t=T(\alpha)\in\alpha^eR[[\alpha]]$ solving
$f(\alpha)=0$; its coefficient of $\alpha^e$ is a unit.
Conversely, in the finite free, $t$-adically complete algebra $B_f$,
the image of $U(t,X)$ is a unit: modulo $t$ its constant term is a
unit and $X$ is nilpotent. Hence $t=-\alpha^e/U(t,\alpha)$ in $B_f$.
The $t$- and $\alpha$-adic topologies are cofinal. The two resulting
continuous maps identify $B_f$ with $R[[\alpha]]$ and prove the
first identity in (D.EA.3). This argument uses no reducedness or
Noetherian hypothesis on $R$.

Substitution of $t=\alpha^e u$ in the derivative shows that each
nonzero term has precisely its weight as an $\alpha$-power, times
a coefficient in $R[[\alpha]]$. There is one least term, of weight
$d$, and its coefficient is a unit. This proves the second identity
in (D.EA.3) over $R$ itself. Thus the derivative becomes a unit after
inverting $t$, proving the asserted etaleness. The standard monogenic
trace determinant is the norm of the derivative, up to sign. The norm
of $\alpha$ is $(-1)^e a_e(t)$, which is $t$ times a unit, and the norm
of a unit is a unit. The discriminant is consequently $t^d$ times
a unit, again over arbitrary $R$.

\emph{The elementary Newton estimate.}
We will use a version which works over every coefficient ring.
In $R[[\alpha]]$, let $h$ be a polynomial and let $\beta$ satisfy
\[
 h'(\beta)=\alpha^d v\quad(v\text{ a unit}),\qquad
 h(\beta)\in\alpha^aR[[\alpha]],\qquad a>2d.
                                                               \tag{D.EA.5}
\]
Division of $h(\beta)$ by $h'(\beta)$ is defined, because
$\alpha$ is a nonzerodivisor. Set
$\beta_1=\beta-h(\beta)/h'(\beta)$. Taylor's polynomial identity
gives
\[
 \beta_1-\beta\in\alpha^{a-d}R[[\alpha]],\qquad
 h(\beta_1)\in\alpha^{2(a-d)}R[[\alpha]].               \tag{D.EA.6}
\]
The derivative changes by an element divisible by $\alpha^{a-d}$,
so its quotient by $\alpha^d$ is still a unit. Iteration is therefore
defined; the excess of the error order over $2d$ at least doubles
at every step. Completeness yields a root $\beta_\infty$, with
$\beta_\infty-\beta\in\alpha^{a-d}R[[\alpha]]$.
This proof uses polynomial Taylor identities, not division by
factorials, and works in every characteristic. It is functorial in
the coefficient ring.

\emph{Coverage of local fields.}
Let $L/k((t))$ have the degree and different in the statement.
Its integer ring is a complete discrete valuation ring with residue
field $k$, so the extension is totally ramified and a uniformizer
$\alpha$ has an Eisenstein minimal polynomial $f_\infty$ of degree
$e$. Its derivative has value $d$ for the valuation normalized by
$v_L(\alpha)=1$. The distinct derivative weights just used show
that its coefficients satisfy the same lowest-weight conditions.
Truncate them modulo $t^N$, obtaining the polynomial $g$ of
(D.EA.2). Terms removed have weights at least $eN>d$, so $g$ lies
in $\mathcal E_{e,d}(k)$. In the original integer ring,
\[
 g(\alpha)\in(t^N)=(\alpha^{eN}),\qquad
 g'(\alpha)=\alpha^d\cdot\text{unit},\qquad eN>2d.
\]
The Newton estimate supplies a root of $g$ in $L$. The polynomial
$g$ is Eisenstein of the same degree $e$, so the corresponding
embedding $k((t))[X]/(g)\hookrightarrow L$ is an isomorphism.
There are $e(N-1)=e\lfloor2d/e\rfloor\le2d$ retained coefficients.
The argument works unchanged after an algebraically closed field
extension. The different/derivative identity and the field-valued
Eisenstein description also appear explicitly in
\citet[Section 2 and equation (4.1)]{Yasuda2025EisensteinB}.

\emph{Finite equations for pointed isomorphisms.}
For two polynomials $f,g$ in the chart, work in the free rank-$e$
algebra
\[
 B_f/t^N B_f=(R[t]/t^N)[\alpha]/(f).
\]
Introduce $eN$ coefficients $b_{ij}$ and the candidate element
\[
 \beta=\sum_{i=0}^{e-1}\sum_{j=0}^{N-1}b_{ij}t^j\alpha^i.
                                                               \tag{D.EA.7}
\]
Impose $b_{00}=0$, make $b_{10}$ invertible by one inverse variable,
and impose the $eN$ coefficient equations
\[
                         g(\beta)=0\pmod{(t^N,f)}.     \tag{D.EA.8}
\]
Use the same polynomial representative (D.EA.7) in $B_f$. Its image
in $R[[\alpha]]$ is $\alpha$ times a unit, since $e\ge2$.
The derivative-weight equations for $g$ therefore give
$g'(\beta)=\alpha^d$ times a unit; this assertion remains true
although $\beta$ is not yet a root. Equations (D.EA.8) give
$g(\beta)\in\alpha^{eN}R[[\alpha]]$. Newton's estimate now gives
an exact root $\beta_\infty$. Since $eN-d>d\ge1$, the correction
does not change its invertible linear coefficient. Thus substitution
$X\mapsto\beta_\infty$ gives an augmentation-preserving
$R[[t]]$-algebra isomorphism $B_g\simeq B_f$: a power series with
zero constant term and invertible linear coefficient is an
automorphism of the pointed formal disc over any $R$.

Conversely, such an isomorphism sends the source uniformizer to an
element $\sum_{i<e}b_i(t)\alpha^i$ of $B_f$. Preservation of the
augmentation says $b_0(0)=0$; invertibility on the conormal module
says $b_1(0)$ is a unit. Truncating every $b_i(t)$ modulo $t^N$
produces (D.EA.7) and (D.EA.8). This proves the equivalence on all
$R$-points. We assert existence, not that each approximate witness
is a unique isomorphism or that its Newton lift retains the entire
$t^N$-jet.

Finally, two coefficient charts use at most $4d+4$ coordinates,
and (D.EA.7) with its inverse uses at most $2d+e+1$. Since
$e\le d+1$, the total is at most $7d+6$. To expand (D.EA.8), the
largest initial $\alpha$-power is $e(e-1)$. Replacing
$\alpha^e$ by $-\sum a_i(t)\alpha^{e-i}$ decreases that power by
at least one and adds at most one to coefficient degree at each
step. Starting with degree at most $e+1$ gives the safe bound
$e^2+1$ for every resulting coefficient equation. All other
equations have degree at most two. The number of equations is
$O(d+e)$, from the displayed coefficient and unit conditions.
\end{proof}

\begin{corollary}[Proved: polynomial-complexity charts for bounded local branch data]
\label{D:local-branch-data-linear-complexity}
Fix a finite ordered collection of nontrivial totally ramified branches
over algebraically closed residue fields, with degrees $e_j\ge2$ and
different exponents $d_j$. Put $D=\sum_jd_j$ and suppose $e_j\le n$.
The product of the coefficient charts above uses at most $2D$
coefficient parameters, together with $O(D)$ inverse variables.
The product pointed-isomorphism relation has an existential witness
using $O(D)$ coordinates and equations of degree at most $n^2+1$.
If $D\le cn$, allowing all such ordered degree/different lists and
all their groupings into branch fibres requires at most
$n^{O_c(n)}$ charts, still with $O_c(n)$ coordinates and polynomial
degree in $n$. Unramified branches of degree one add no parameters.
Every counted branch fibre is required to contain a nontrivial branch;
arbitrarily many additional all-unramified marked fibres are not included
in this combinatorial count.
These are charts for local branch data, not a Hurwitz atlas for covers
of a proper curve and not a bound for covers of a surface puncture.
\end{corollary}
\begin{proof}
Add the coordinate bounds in the proposition. Since $d_j\ge e_j-1$,
the number of nontrivial branches is at most $D$ and
$\sum e_j\le2D$. An ordered list of positive integers $d_j$ with
sum at most $D$ has at most $2^{D+1}$ possibilities. For each such
list, choose $2\le e_j\le d_j+1$; a crude upper bound for the
number of choices is $(D+1)^D$. Placing separators to group a list
into branch fibres contributes at most $2^D$. Trivial branches
in each fibre, if total covering degree $n$ is specified, are
then determined by subtracting its nontrivial degrees from $n$.
The same bounds with an additional factor $(n+1)^{D+1}$ allow an
arbitrary bounded total degree instead. These elementary bounds
are $n^{O_c(n)}$ when $D\le cn$; the cases $n=1$ or $D=0$ are
trivial and may be treated separately. Products of the witness
equations retain their degree bound. No global gluing or effective
algebraization statement is needed for these local assertions.
\end{proof}

\begin{remark}[The remaining global step]
To count finite covers of one fixed proper curve, the local data must
still be glued to a cover of its complement. The preceding charts
do not bound the degree of that gluing problem or the number of
global covers having specified formal local extensions. Likewise,
the vanishing of a local deformation obstruction would not by
itself bound the number of connected components of a global
Hurwitz space. Those are additional quantitative gates, retained
explicitly here.
\end{remark}
